\section{Two outer splits and five residual witnesses}\label{sec:outer}

We first transmit the finite computation through prediction trials whose
boundaries are selected by markers. When a required boundary is absent,
a second split uses a long periodic interval to keep the prefix metadata
constant while varying the clock residue. Each construction defines a total
affine update before describing its successful cuts. The episodes where
neither construction supplies every input with a cut will have one of five
explicit witnesses.

The schedules in this section have fixed finite lengths and offsets for
the chosen monoid and alphabet. We state every separation and visibility
condition when it is used. Proposition~\ref{prop:parameters} constructs
all of them in one consistent order; until then they are the explicit
hypotheses of the two split constructions.

Use the episode start \(s\), sampling origin \(t=s+L_0\), anchor \(q\), and end
\(q+B_{\rm tail}\) from Definition~\ref{def:episodes}. Five later subbuffers
will lie between \(s\) and \(t\); the two splits of this section make their cuts
after \(t\). Write \(T_e=T_{s,q+B_{\rm tail}}\).

\subsection{The main schedule and the common end guard}

Fix the initial representation on \(V\), and put \(J=|V|^2\).
There is a finite collection of \emph{main trials}, indexed by \(\nu\),
with fixed predictions \(\Theta_\nu\in M^{J+1}\).
Let
\[
             \mathcal Z=\{(\nu,j):0\le j\le J+1\}
\]
label their endpoints and checkpoints. For \(\lambda\in\mathcal Z\), let
\(o_\lambda\) be its nominal offset from \(t\). Use an inner marker system
of separation \(d_\circ\), block length \(K_\circ>3d_\circ^2\), and radius
\(r_\circ\ge K_\circ\). Define
\begin{equation}\label{eq:main-boundaries}
\begin{split}
 z_\lambda(t)&=t+o_\lambda,\\
 u_\lambda(t)&=\text{first inner marker in }
 [z_\lambda(t)-H,z_\lambda(t)+H],
\end{split}
\end{equation}
with value \(\bot\) when the search is empty.
Let \(\mathcal A(t)\) mean that every \(u_\lambda(t)\) is present.
On \(\mathcal A(t)\), these are the actual ordered trial boundaries.
For trial \(\nu\), its actual endpoints are
\(B_\nu=u_{\nu,0}(t)\), \(C_\nu=u_{\nu,J+1}(t)\).
Its block \(j\) consists of cuts
\begin{equation}\label{eq:main-cuts}
                 k=t+c_{\nu,j}+a,\qquad 0\le a<R_1.
\end{equation}
The offsets place this entire block between boundaries \(j-1\) and \(j\).
All earlier-boundary searches, including their inspection neighborhoods, are
before the block; all later ones are after it. Trials are disjoint and ordered.

The second layer will use cuts
\begin{equation}\label{eq:periodic-cuts}
                     k=z_\lambda(t)+a,\qquad 0\le a<R_2.
\end{equation}
Let \(\mathcal H_1,\mathcal H_2\) be the respective full finite sets of cut
offsets from \(t\) in \eqref{eq:main-cuts} and \eqref{eq:periodic-cuts}.
Choose the end-stability radius with
\[
              \diam\mathcal H_1,\diam\mathcal H_2<w_E<d_\bullet,
 \qquad K_\bullet-w_E>d_\bullet^2.
\]

Every suffix test in outer layer \(r\in\{1,2\}\) uses the following guard,
independently of its tag. At its beginning \(k\), form \emph{all} origins
\[
                         k-h,\qquad h\in\mathcal H_r.
\]
Require their values of \(\eta\) to be equal, and require exact consumption
to the episode end obtained from one fixed reference origin in this set.
Let \(q\) be that reference anchor. All bounded guard data must be present
in the suffix argument.

If paired with an actual prefix cut, the true \(t\) belongs to this list.
Lemma~\ref{lem:shared-anchor} forces the reference anchor to be the true \(q\),
even if the proposed suffix initially extends beyond the first episode or
ends prematurely. Thus it consumes exactly the first episode. In particular
the clock integers supplied on the two sides,
\[
                            k-t,\qquad q-k,
\]
always have common total
\begin{equation}\label{eq:clock-total}
                              \ell=q-t.
\end{equation}
For availability, it suffices that \(\eta\) is stable under every shift of
magnitude at most \(w_E\), since this covers every reference origin relative
to the actual one.

\subsection{The main affine update}

Apply Lemma~\ref{lem:clock} to tags
\[
                 \alpha=(\nu,j,X,Y),\qquad X,Y\in V.
\]
Write the resulting clock as \(n_1,I^{(1)}_\alpha,J^{(1)}_\alpha,W_1\),
and let \(\mathcal G_1(\ell)\) be its graph of possible tag pairs.
Take \(R_1=n_1\).

Define \(F_{1,e}:V\to V\) for every episode, without referring to any split.
\begin{enumerate}[label=\textup{(\alph*)}]
\item If \(\mathcal G_1(\ell)\) is contained in the diagonal and
\(\mathcal A(t)\) holds, process letters normally except that each main trial
has the ideal deterministic shift of Lemma~\ref{lem:prediction} added at its end.
\item If \(\mathcal G_1(\ell)\) is contained in the diagonal and
\(\mathcal A(t)\) fails, use \(v\mapsto vT_e\).
\item If \(\mathcal G_1(\ell)\) has its unique off-diagonal edge
\(\alpha\to\beta\), use
\begin{equation}\label{eq:main-error}
                 F_{1,e}(v)=vT_e+Y(\beta)-X(\alpha)T_e.
\end{equation}
\end{enumerate}
These are exhaustive cases, including the empty graph in the diagonal case.
In every case \(F_{1,e}\) is deterministic affine with linear part \(T_e\).

Here and below a language with metadata is the finite union over
\emph{its own} metadata choices. Prefix and suffix unions are independent.
Let \(P_x^{(1)}\) require, for one tag \((\nu,j,X,Y)\):
\begin{enumerate}[label=\textup{(\roman*)}]
\item its length is one prescribed cut in \eqref{eq:main-cuts};
\(k-t\in I^{(1)}_\alpha\), and \(x=X\);
\item every boundary preceding the cut is present, including the boundaries
of all earlier trials and through boundary \(j-1\) of the current trial;
\item processing from \(x\) to \(B_\nu\), with ideal shifts on earlier trials,
arrives at the enumerated input \(x_j\) of this block;
\item the block's predicted transfer and all its preceding individual
interval conditions in Definition~\ref{def:eligibility} hold.
\end{enumerate}
All these data are before the cut, by placement.

Let \(Q_y^{(1)}\) use its own tag \((\nu,j,X,Y)\), the independent end guard,
the residue check \(q-k\in J^{(1)}_\alpha\), and \(y=Y\).
For each candidate
\begin{equation}\label{eq:main-role-origins}
                    k-c_{\nu,j}-a,\qquad 0\le a<R_1,
\end{equation}
perform the future-boundary searches for that role. Require every boundary
from \(j\) onward in the current trial, and every boundary in later trials,
to be present at the \emph{same word position} for all these candidates.
These searches are after \(k\) for every fixed role, whether the role is true
or false. Use their common positions to check the suffix-product condition
from boundary \(j\) to \(C_\nu\). From the block's enumerated output \(y_j\)
at \(C_\nu\), compute onward with ideal shifts in later trials and require
final result \(y\).

Define \(S_1\subseteq E\) by the following conditions on actual episode data:
\begin{equation}\label{eq:S1}
\begin{gathered}
 \eta(t)\text{ is stable within }w_E,\qquad
 \ell\notin W_1,\qquad \mathcal A(t),\\
 u_\lambda(t)\text{ is stable within }R_1\text{ for every }\lambda,\qquad
 \text{some main trial is perfectly predicted}.
\end{gathered}
\end{equation}
Stability means equality of the returned position or of the absent value for
every integer shift of either sign up to the indicated bound.

\begin{lemma}[Main split]\label{lem:main-split}
The languages \(P_x^{(1)},Q_y^{(1)}\) satisfy
Lemma~\ref{lem:split} on \(D=E\), with updates \(F_{1,e}\) and skipped class
\(E\setminus S_1\).
\end{lemma}

\begin{proof}
Take an arbitrary accepted prefix/suffix pair at a genuine episode boundary.
The end guard gives exact consumption and the common total \(\ell\).
If \(\mathcal G_1(\ell)\) is diagonal, the tags match. The actual \(t\)
then appears among \eqref{eq:main-role-origins}. The prefix establishes all
earlier boundaries and the suffix establishes all later boundaries, so
\(\mathcal A(t)\) holds. Their joint checks give an ideally eligible pair of
the actual current trial. Lemma~\ref{lem:prediction}, together with the ideal
computations before and after that trial, gives \(F_{1,e}(x)=y\).

If tags differ, the clock graph consists of that unique off-diagonal edge.
The input and output checks fix \(x=X(\alpha)\), \(y=Y(\beta)\), so
\eqref{eq:main-error} fits the transfer. The calculations performed under
possibly false roles introduce no competing input or output. These cases prove
soundness for every independent pair of witnesses.

For \(e\in S_1\) and input \(x\), select a perfect trial and its block for the
vector actually arriving there and its trial output. Choose the tag
\(X=x\), \(Y=F_{1,e}(x)\). Since \(\ell\notin W_1\), this diagonal tag pair
is available at some first residue. Its \(n_1\) consecutive cut choices realize
every such residue. Every role-origin differs from the true one by less than
\(R_1\), so future-boundary stability makes all suffix comparisons pass.
End stability makes the independent guard pass. All required data are present,
giving the desired split.
\end{proof}

The first split is sound on every episode, but its availability needs all
main boundaries and a perfect prediction. We next handle an episode with
a sufficiently long markerless boundary window. Its short-period run will
supply many cuts with identical prefix products and local contexts; only
the second clock residue changes.

\subsection{The periodic-window affine update}

To compute \(F_1\) on skipped sequences, use its homogeneous lift
\(\widehat F_{1,e}\) on
\[
                         U_0=V\oplus\mathbf F_2.
\]
Take a second clock with modulus \(n_2\), bad set \(W_2\), and tags
\begin{equation}\label{eq:second-tags}
 \alpha=(\lambda,X,Y,m,r,\xi),\qquad
 \begin{cases}
 \lambda\in\mathcal Z,\quad X,Y\in U_0,\\
 m\in M,\quad r\in\mathbb Z/n_1,\\
 \xi\in\Sigma^{2r_\circ}.
 \end{cases}
\end{equation}
These fields propose the boundary label, episode input/output, product
\(T_{s,k}\), residue \(k-t\bmod n_1\), and immediate left context of the cut.
There is no field for the cut index or \(n_2\).

Let \(\mathcal G_2(\ell)\) be its graph. On every episode define \(F_{2,e}\) to
be \(\widehat F_{1,e}\) in the diagonal case; in the unique off-diagonal case
\(\alpha\to\beta\), define
\begin{equation}\label{eq:second-error}
 F_{2,e}(v)=v\widehat F_{1,e}
       +Y(\beta)-X(\alpha)\widehat F_{1,e}.
\end{equation}
This is deterministic affine with linear part \(\widehat F_{1,e}\).

The prefix language \(P_x^{(2)}\) chooses its own tag and requires:
\begin{enumerate}[label=\textup{(\roman*)}]
\item a cut in \eqref{eq:periodic-cuts} for its label,
\(k-t\in I^{(2)}_\alpha\), and \(x=X\);
\item the proposed product \(m=T_{s,k}\), first-clock residue \(r\), and context
\(\xi\) are correct;
\item there is no inner marker in
\[
                           [z_\lambda(t)-H,k-r_\circ].
\]
\end{enumerate}
The marker decisions use only prefix data, since their radius is \(r_\circ\).

The suffix language \(Q_y^{(2)}\) uses its own tag, the second layer's independent
end guard, \(q-k\in J^{(2)}_\alpha\), and \(y=Y\).
Adjoin the proposed \(\xi\) immediately before its argument for the following
local test: require no inner marker in
\begin{equation}\label{eq:periodic-suffix-window}
                          [k-r_\circ,k+H+R_2].
\end{equation}
Exactly \(2r_\circ\) preceding letters suffice for the earliest radius
\(r_\circ\) decision. All other required letters must be in the suffix.

Use the proposed \(m\) and the actual suffix product to obtain a proposed
\(T_e\). Use \(r+(q-k)\bmod n_1\) as the proposed first-clock total.
Under the hypothesis \(\neg\mathcal A(t)\), these two data determine \(F_1\):
it is the linear action in the diagonal first-clock case and
\eqref{eq:main-error} in the off-diagonal case. Require
\[
                             y=X\widehat F_{1,e}
\]
as evaluated by precisely this rule. The suffix need not decode \(t\).

Put \(N_M=|M|!\), and require
\begin{equation}\label{eq:periodic-parameters}
\begin{gathered}
 R_2=n_2n_1N_Md_\circ,\qquad
 H-d_\circ>2r_\circ+|M|d_\circ+R_2,\\
 \text{every prime factor of }n_2\text{ exceeds }d_\circ,N_M,n_1.
\end{gathered}
\end{equation}
Within \(E\setminus S_1\), let \(S_2\) be the episodes with
\begin{equation}\label{eq:S2}
\begin{gathered}
 \eta(t)\text{ stable within }w_E,\qquad \ell\notin W_2,\\
 \text{for some }\lambda,\quad
 [z_\lambda(t)-H,z_\lambda(t)+H+2R_2]\text{ has no inner marker}.
\end{gathered}
\end{equation}

\begin{lemma}[Periodic-window split]\label{lem:periodic-split}
The second pair of languages satisfies Lemma~\ref{lem:split} on
\(D=E\setminus S_1\), with updates \(F_{2,e}\) and skipped class
\[
                        D_0=E\setminus(S_1\cup S_2).
\]
\end{lemma}

\begin{proof}
Again the independent guard gives exact consumption and total \(\ell\).
For matching tags, their product, residue and context are actual.
Since \(k=z_\lambda(t)+a\), \(0\le a<R_2\), the prefix absence interval and
\eqref{eq:periodic-suffix-window} cover the full actual search
\([z_\lambda(t)-H,z_\lambda(t)+H]\). Thus this boundary is absent and
\(\neg\mathcal A(t)\) is proved by the successful split.
The suffix has consequently evaluated \(\widehat F_{1,e}\) correctly.
In the diagonal second-clock case its output is \(F_{2,e}(x)\).
Mismatched tags fit the unique input/output pair in \eqref{eq:second-error}.
This proves soundness without assuming correctness of unmatched context.

For availability, let \(e\in S_2\), with the indicated center \(z\).
The marker property gives short-period \(K_\circ\)-blocks at all starts
from \(z-H+d_\circ\) through \(z+H-d_\circ\).
Lemma~\ref{lem:periodic-overlap} makes these one periodic run of some
period \(d<d_\circ\). The second inequality of
\eqref{eq:periodic-parameters} places all needed contexts and at least
\(|M|\) whole periods immediately before \(z\) within this run.

For any \(c\in M\), two of \(c^0,\ldots,c^{|M|}\) agree. Multiplying onward
gives an eventual period at most \(|M|\), starting before \(|M|\), and that
period divides \(N_M\). In particular
\[
                    c^{r+N_M}=c^r\qquad(r\ge|M|).
\]
Consider cuts
\begin{equation}\label{eq:special-cuts}
                k_a=z+a\,dN_Mn_1,\qquad 0\le a<n_2.
\end{equation}
They have \(0\le k_a-z<R_2\), identical \(2r_\circ\)-letter contexts,
and a fixed residue modulo \(n_1\). Take the period word in phase at \(z\),
with product \(c=T_{z,z+d}\), and put \(A=T_{s,z-|M|d}\). Then
\[
 T_{s,k_a}=Ac^{|M|+aN_Mn_1}=Ac^{|M|}.
\]
Thus the prefix-product field \(m\) is fixed. Although the suffix product
can vary with \(a\), it always satisfies \(mT_{k_a,q+B_{\rm tail}}=T_e\);
the recovered first-clock total is likewise always \(q-t\).
For input \(x\), one correct tag in \eqref{eq:second-tags}, with
\(X=x\) and \(Y=x\widehat F_{1,e}\), is therefore valid at all these cuts.
The prime-factor requirement gives
\(\gcd(dN_Mn_1,n_2)=1\). The cuts realize every residue modulo \(n_2\),
so one realizes this diagonal tag pair since \(\ell\notin W_2\).
Both absence tests lie in the markerless interval from \eqref{eq:S2}:
their rightmost endpoint is less than \(z+H+2R_2\), and their leftmost
positions are covered by the same margin. The stable end guard passes.
This gives a split for every input.
\end{proof}

We have two total updates and two availability sets. The remaining task
is to describe failure of both availability conditions using a fixed
number of predicates. These predicates will locate the true cut in the
aligned constructions of Section~\ref{sec:aligned}.

\subsection{The five witnesses}

We use the following predicates \(A_i(t,q)\) on episode data, and also at
formally translated sampling origins:
\begin{equation}\label{eq:witnesses}
\begin{array}{c@{\quad}p{0.80\linewidth}}
 i & \(A_i(t,q)\)\\ \hline
 0 & \(q-t\in W_1\), or \(\mathcal A(t)\) and no main trial is perfectly
 predicted at these boundaries;\\
 1 & some \(u_\lambda(t)\) changes under a shift of magnitude at most \(R_1\);\\
 2 & \(q-t\in W_2\);\\
 3 & some \([z_\lambda(t)-H,z_\lambda(t)+H]\) is markerless, but an inner
 marker lies among the next \(2R_2\) positions to its right;\\
 4 & \(\eta(t)\) changes under a shift of magnitude at most \(w_E\).
\end{array}
\end{equation}
In particular \(A_4\) does not use \(q\). Membership in a clock bad set always
uses the corresponding modulus.

At a translated origin, these are \emph{formal predicates on the given data},
with \(q\) fixed. They do not assert that the origin completes another episode,
and do not run another first-break scan. All marker decisions and products
they request will be internal to the available data.

\begin{lemma}[Residual coverage]\label{lem:witness-coverage}
Every episode of \(D_0\) satisfies at least one of \(A_0,\ldots,A_4\).
\end{lemma}

\begin{proof}
Suppose \(A_4\) fails. If \(\mathcal A(t)\) holds and both \(A_0,A_1\) fail,
all conditions of \(S_1\) hold. If \(\mathcal A(t)\) fails, choose a missing
boundary. Failure of \(A_2,A_3\) gives a good second-clock total and an empty
window extended by \(2R_2\) at that boundary, hence \(S_2\) unless the episode
was already in \(S_1\). Thus an episode in neither designation has a witness.
\end{proof}
